\documentclass[11pt]{beamer}

\usetheme{Madrid}
\usecolortheme{default}

\usepackage[spanish]{babel}
\usepackage[utf8]{inputenc}
\usepackage[T1]{fontenc}
\usepackage{lmodern}

\usepackage{amsmath,amssymb,amsfonts}
\usepackage{graphicx}
\usepackage{tikz}
\usetikzlibrary{arrows.meta,positioning}

\usepackage[backend=biber]{biblatex}
\addbibresource{./bibliografia.bib}

\newcommand{\card}{\operatorname{Card}}
\newcommand{\floor}[1]{\lfloor #1 \rfloor}
\newcommand{\ceil}[1]{\lceil #1 \rceil}
\newcommand{\partfrac}[1]{\langle #1 \rangle}

\AtBeginSubsection[]
{
\begin{frame}{Índice de contenidos}
\tableofcontents[
currentsection,
currentsubsection
]
\end{frame}
}

\AtBeginSection[]
{
\begin{frame}{Índice de contenidos}
\tableofcontents[
currentsection,
currentsubsection
]
\end{frame}
}

\newenvironment{demostracion}
{\par\noindent\textbf{Dem.}\ }

% Colores
% \definecolor{UMGranate}{RGB}{120,0,40}

% \setbeamercolor{structure}{fg=UMGranate}
% \setbeamercolor{frametitle}{bg=UMGranate,fg=white}

% \setbeamercolor{block title}{bg=UMGranate,fg=white}
% \setbeamercolor{block body}{bg=gray!10}

% \setbeamercolor{block title alerted}{bg=UMGranate,fg=white}
% \setbeamercolor{block body alerted}{bg=gray!10}

\title[Trabajo de Fin de Grado]
{\Large El teorema de equidistribución de Weyl \\ [0.1cm] 
\small Trabajo de Fin de Grado}

\author[D. A. Manzanares Sánchez]
{\textbf{Domingo Antonio Manzanares Sánchez} \\ [0.3cm]
Tutor: Gustavo Adolfo Garrigós Aniorte}

\institute[UM]{Facultad de Matemáticas \\
Universidad de Murcia}

\date{Junio de 2026}

\begin{document}

\begin{frame}
\titlepage
\end{frame}

\begin{frame}
\frametitle{Índice de contenidos}
\tableofcontents
\end{frame}

\section{Introducción}

\begin{frame}{Introducción}

Una de las sucesiones más estudiadas en teoría de números es
\[
  \partfrac{\alpha}, \, 
  \partfrac{2\alpha}, \, 
  \partfrac{3\alpha}, \, 
  \ldots, \, 
  \partfrac{n\alpha}, \, 
  \ldots
\]
donde $\partfrac{ x }$ denota la parte fraccionaria de $x$.

\vspace{0.3cm}

\begin{itemize}
\item Si $\alpha$ es racional, la sucesión es periódica.
\item Si $\alpha$ es irracional, la sucesión tiene infinitos términos distintos.
\end{itemize}

\vspace{0.3cm}
\pause

\begin{block}{Pregunta}
¿Cómo se distribuyen estos valores sobre el intervalo $[0,1)$?
\end{block}

\end{frame}

\section{Equidistribución módulo 1}



\begin{frame}
\frametitle{Definición de W}
\end{frame}

\begin{frame}{Números reales módulo 1}

\begin{block}{Parte fraccionaria}
Para $x\in\mathbb R$ definimos
\[\langle x\rangle=x-\lfloor x\rfloor.\]
Por tanto,
$\langle x\rangle\in[0,1).$
\end{block}
\pause
\begin{block}{Congruencia módulo 1}
Dos números reales son equivalentes si difieren en un entero:
\[
x\equiv y \pmod 1
\quad\Longleftrightarrow\quad
x-y\in\mathbb Z.
\]
\end{block}
\pause
Cada clase de equivalencia tiene un representante único en $[0,1)$.
\end{frame}

\subsection{AAAAAAAAAAA}

\begin{frame}{El toro unidimensional}
\[
\mathbb T=\mathbb R/\mathbb Z
\]
\vspace{0.4cm}

\begin{itemize}
\item Identificamos los puntos $0$ y $1$.
\item El intervalo $[0,1)$ se convierte en una circunferencia.
\item Las funciones en $\mathbb T$ son exactamente las funciones
      periódicas de periodo $1$.
\end{itemize}

\vspace{0.3cm}

\begin{center}
\Large
$0 \equiv 1$
\end{center}
\end{frame}

\begin{frame}{Equidistribución módulo 1}

\begin{block}{Definición}
Una sucesión $(x_n)$ es equidistribuida módulo $1$ si para todo intervalo
$(a,b)\subset [0,1)$,

\[
\lim_{N\to\infty}
\frac{A_N((a,b),x)}{N}
=
b-a.
\]
\end{block}

\pause

donde

\[
A_N((a,b),x)
=
\#\{1\le n\le N:\langle x_n\rangle\in(a,b)\}.
\]

\pause

\begin{alertblock}{Interpretación}
La proporción de términos que cae en un intervalo coincide con
la longitud del intervalo.
\end{alertblock}

\end{frame}

\begin{frame}{Ejemplos}

\textbf{No equidistribuida}

\[
0,\;0,\;\frac12,\;0,\;\frac13,\;\frac23,\ldots
\]

Los términos se concentran excesivamente cerca de ciertos puntos.

\pause

\vspace{0.4cm}

\textbf{Equidistribuida}

\[
0,\;
0,\frac12,\;
0,\frac13,\frac23,\;
0,\frac14,\frac24,\frac34,\ldots
\]

La distribución se vuelve cada vez más uniforme.

\pause

\vspace{0.4cm}

\textbf{Idea clave}

\[
\text{Densidad}
\neq
\text{Equidistribución}.
\]

\end{frame}

\begin{frame}{La sucesión $n\alpha$}

\[
\langle \alpha\rangle,
\langle 2\alpha\rangle,
\langle 3\alpha\rangle,
\dots
\]

\pause

\begin{block}{Caso racional}
Si

\[
\alpha=\frac pq,
\]

la sucesión es periódica y toma un número finito de valores.
\end{block}

\pause

\begin{block}{Caso irracional}
Si

\[
\alpha\notin\mathbb Q,
\]

los valores nunca se repiten y se distribuyen uniformemente en
$\mathbb T$.
\end{block}

\end{frame}

\begin{frame}{Teorema de Weyl}
\begin{alertblock}{Teorema de Weyl}
Si $\alpha$ es irracional, entonces

\[
(n\alpha)_{n\ge1}
\]

es equidistribuida módulo $1$.
\end{alertblock}

\end{frame}

\begin{frame}{Números reales módulo 1}

\begin{block}{Parte fraccionaria}
\[
\langle x\rangle = x-\lfloor x\rfloor,
\qquad
\langle x\rangle\in[0,1).
\]
\end{block}

\vspace{0.2cm}

\begin{block}{Relación de equivalencia}
\[
x\equiv y \pmod 1
\quad \Longleftrightarrow \quad
x-y\in\mathbb Z.
\]
\end{block}

\vspace{0.2cm}

\begin{exampleblock}{Ejemplos}
\[
0.3 \equiv 1.3 \equiv 2.3 \pmod 1
\]

\[
-0.7 \equiv 0.3 \pmod 1
\]
\end{exampleblock}

\end{frame}

\begin{frame}{El toro unidimensional}

\begin{block}{Espacio cociente}
\[
\mathbb T=\mathbb R/\mathbb Z.
\]
\end{block}

\begin{itemize}
\item Cada clase tiene un único representante en $[0,1)$.
\item Identificamos los extremos:
\[
0 \equiv 1.
\]
\item La recta real se enrolla sobre la circunferencia unidad.
\end{itemize}

\vspace{0.3cm}

\begin{center}
% Dibujo o imagen del toro
% \includegraphics[width=.6\textwidth]{toro.pdf}
\end{center}

\end{frame}

\begin{frame}{Funciones sobre $\mathbb T$}

\begin{block}{Observación}
Una función sobre $\mathbb T$ es equivalente a una función
$1$-periódica sobre $\mathbb R$.
\end{block}

\[
f(x+1)=f(x).
\]

\vspace{0.3cm}

\[
f(x)=f(\langle x\rangle).
\]

\vspace{0.3cm}

\begin{center}
% \includegraphics[width=.7\textwidth]{periodica.pdf}
\end{center}

\end{frame}

\begin{frame}{Equidistribución módulo 1}

\begin{block}{Definición}
Una sucesión $(x_n)$ es equidistribuida módulo $1$ si

\[
\lim_{N\to\infty}
\frac{A_N((a,b),x)}{N}
=
b-a
\]

para todo intervalo $(a,b)\subset\mathbb T$.
\end{block}

\vspace{0.2cm}

\[
A_N((a,b),x)
=
\#\{1\le n\le N:\langle x_n\rangle\in(a,b)\}.
\]

\vspace{0.3cm}

\begin{alertblock}{Interpretación}
La proporción de términos que cae en un intervalo coincide
con la longitud del intervalo.
\end{alertblock}

\end{frame}

\begin{frame}{Propiedades y ejemplos}

\begin{columns}

\column{0.48\textwidth}

\begin{block}{Propiedades}
\begin{itemize}
\item Eliminar un número finito de términos.
\item Invarianza por traslaciones.
\item Intercalación de sucesiones equidistribuidas.
\end{itemize}
\end{block}

\column{0.48\textwidth}

\begin{exampleblock}{Ejemplos}
\begin{itemize}
\item Sucesiones equidistribuidas.
\item Sucesiones no equidistribuidas.
\end{itemize}
\end{exampleblock}

\end{columns}

\vspace{0.3cm}

\begin{center}
% Dibujos de puntos sobre [0,1)
% \includegraphics[width=.8\textwidth]{ejemplos.pdf}
\end{center}

\end{frame}

\begin{frame}{La sucesión $n\alpha$}

\begin{block}{Proposición}
Sea $\alpha\in\mathbb R$.
\begin{itemize}
\item Si $\alpha\in\mathbb Q$, la sucesión
$(\langle n\alpha\rangle)$ es periódica.
\item Si $\alpha\notin\mathbb Q$, posee infinitos términos distintos.
\end{itemize}
\end{block}

\vspace{0.2cm}

\begin{columns}

\column{0.48\textwidth}

\centering
\textbf{$\alpha=\frac13$}

\vspace{0.2cm}

% \includegraphics[width=\textwidth]{alpha_racional.pdf}

\column{0.48\textwidth}

\centering
\textbf{$\alpha=\sqrt2$}

\vspace{0.2cm}

% \includegraphics[width=\textwidth]{alpha_irracional.pdf}

\end{columns}

\end{frame}

\begin{frame}{Teorema de Weyl}

\begin{alertblock}{Teorema de Weyl}
Si $\alpha\notin\mathbb Q$, entonces

\[
(\langle n\alpha\rangle)_{n\ge1}
\]

es equidistribuida módulo $1$.
\end{alertblock}

\vspace{0.3cm}

\begin{block}{Idea de la demostración}
\begin{enumerate}
\item Se prueba el resultado para exponenciales
      $e^{2\pi ikx}$.
\item Se extiende a polinomios trigonométricos.
\item Se aproxima la función característica mediante
      funciones continuas.
\end{enumerate}
\end{block}

\vspace{0.2cm}

\begin{center}
% \includegraphics[width=.55\textwidth]{aproximacion.pdf}
\end{center}

\end{frame}
\section{Caracterizaciones de la equidistribución}
\begin{frame}{¿Cómo verificar la equidistribución?}

\begin{block}{Problema}
La definición exige comprobar que
\[
\lim_{N\to\infty}\frac{A_N((a,b),x)}{N}=b-a
\]
para todo intervalo $(a,b)\subset\mathbb T$.
\end{block}

\pause

\begin{alertblock}{Objetivo}
Buscar criterios más prácticos para demostrar
que una sucesión es equidistribuida.
\end{alertblock}

\vspace{0.3cm}

En este capítulo presentamos dos herramientas:
\begin{enumerate}
\item Caracterización funcional.
\item Criterio de Weyl.
\end{enumerate}

\end{frame}

% Caracterización funcional
\begin{frame}{Caracterización funcional}

\begin{block}{Teorema}
Las siguientes afirmaciones son equivalentes:
\begin{enumerate}
\item $(x_n)$ es equidistribuida módulo $1$.
\item Para toda función Riemann integrable
$f:\mathbb T\to\mathbb R$,
\[
\lim_{N\to\infty}\frac{1}{N}\sum_{n=1}^{N}f(x_n)=\int_0^1 f(x)\,dx.
\]
\end{enumerate}
\end{block}

\begin{alertblock}{Interpretación}
La media discreta sobre la sucesión coincide
con la media continua sobre el toro.
\end{alertblock}

\end{frame}

% Idea de geometrica
\begin{frame}{Interpretación geométrica}

\[
\frac1N\sum_{n=1}^{N}f(x_n)
\]

aproxima

\[
\int_0^1 f(x)\,dx.
\]

\vspace{0.4cm}

\begin{center}
% \includegraphics[width=.7\textwidth]{media_integral.pdf}
\end{center}

\vspace{0.2cm}

\begin{alertblock}{Idea}
Una sucesión equidistribuida se comporta como una
muestra uniforme del intervalo $[0,1)$.
\end{alertblock}

\end{frame}

% Criterio de Weyl
\begin{frame}{Criterio de Weyl}

\begin{block}{Teorema}
Las siguientes afirmaciones son equivalentes:

\begin{enumerate}
\item $(x_n)$ es equidistribuida módulo $1$.
\item Para todo $k\in\mathbb Z\setminus\{0\}$,

\[
\frac1N
\sum_{n=1}^{N}
e^{2\pi i k x_n}
\longrightarrow 0.
\]

\end{enumerate}
\end{block}

\pause

\begin{alertblock}{Ventaja}
El problema se reduce al estudio de sumas
de exponenciales complejas.
\end{alertblock}

\pause

\[
\text{Equidistribución}
\quad \Longrightarrow \quad
\text{Análisis de Fourier}
\]

\end{frame}

% Aplicación del criterio de Weyl
\begin{frame}{Aplicación: la sucesión $n\alpha$}

\begin{block}{Resultado}
Si $\alpha\notin\mathbb Q$, entonces $(n\alpha)_{n\geq1}$
es equidistribuida módulo $1$.
\end{block}

Por el criterio de Weyl basta estudiar
\[\frac{1}{N}\sum_{n=1}^{N}e^{2\pi i k n\alpha}.\]

Esta suma es geométrica y satisface
\[
\left|\frac{1}{N}\sum_{n=1}^{N} e^{2\pi i k n\alpha}\right|
\le\frac{2}{N|e^{2\pi i k\alpha}-1|}.
\]

Por tanto,
\[\lim_{N\to\infty}\frac{1}{N}\sum_{n=1}^{N}e^{2\pi i k n\alpha}=0.\]
\end{frame}
\section{Método de Euler}
\begin{frame}{Método de Euler}

\begin{block}{Problema}
Por el criterio de Weyl, para demostrar que una sucesión es
equidistribuida debemos estudiar sumas de la forma
\[\sum_{n=1}^{N} e^{2\pi i f(n)}.\]
\end{block}

\begin{alertblock}{Idea}
Transformar sumas discretas en integrales que puedan
estimarse con precisión.
\end{alertblock}
\vspace{0.3cm}
Herramienta principal:
\[
\boxed{\text{Fórmula de sumación de Euler}}
\]
\end{frame}

\begin{frame}{Fórmula de sumación de Euler}

\begin{block}{Proposición}
Sea $f:[1,N]\to\mathbb C$ con $f'\in C([1,N])$.

Entonces
\[
\sum_{n=1}^{N}f(n)=
\int_1^N f(x)\,dx + \frac{f(1)+f(N)}{2} +
\int_1^N\left(\langle x\rangle-\frac12\right)f'(x)\,dx.
\]
\end{block}

\begin{alertblock}{Interpretación}
La suma discreta se expresa mediante integrales
más fáciles de estimar.
\end{alertblock}

\end{frame}

\begin{frame}{Ejemplo: la sucesión $a\log n$}

\begin{block}{Resultado}
La sucesión $(a\log n)$ no es equidistribuida módulo $1$.
\end{block}

Aplicando la fórmula de Euler a
\[
f(x)=e^{2\pi i a\log x}
\]
obtenemos una estimación de la suma exponencial asociada.

\begin{alertblock}{Conclusión}
El criterio de Weyl falla para $k=1$.
Por tanto, $(a\log n)$ no es equidistribuida.
\end{alertblock}

\end{frame}

\begin{frame}{Ejemplo: la sucesión $an^\sigma$}

\begin{block}{Hipótesis}
\[
a\neq 0,
\qquad
0<\sigma<1.
\]
\end{block}

Aplicando nuevamente la fórmula de Euler a

\[
f(x)=e^{2\pi i a x^\sigma}
\]

se obtiene

\[
\sum_{n=1}^{N}
e^{2\pi i a n^\sigma}
=
o(N).
\]


\begin{alertblock}{Por el criterio de Weyl}
\[
(an^\sigma)
\]

es equidistribuida módulo $1$.
\end{alertblock}

\end{frame}

\begin{frame}{Hacia el teorema de Fejér}

Los ejemplos anteriores sugieren que la equidistribución
depende del comportamiento de la derivada.

\vspace{0.3cm}

\begin{center}

\[
g(x)=\log x
\]

\[
g(x)=x^\sigma
\]

\end{center}

\vspace{0.3cm}

\pause

\begin{alertblock}{Pregunta}
¿Qué condiciones sobre una función $g$
garantizan que $(g(n))$ sea equidistribuida?
\end{alertblock}

\end{frame}


\begin{frame}{Teorema de Fejér}

\begin{block}{Teorema}
Sea $g:[1,\infty)\to\mathbb R$ derivable.

Si

\[
g'(x)\to 0,
\]

\[
x|g'(x)|\to\infty,
\]

y $g'(x)$ es monótona para $x$ suficientemente grande,

entonces

\[
(g(n))
\]

es equidistribuida módulo $1$.
\end{block}

\end{frame}

\begin{frame}{Aplicaciones del teorema de Fejér}

\begin{block}{Consecuencias}
Las siguientes sucesiones son equidistribuidas módulo $1$:
\end{block}

\[
(an^\sigma),
\qquad
0<\sigma<1.
\]

\[
(an^\sigma \log^\tau n).
\]

\[
(a\log^\tau n),
\qquad
\tau>1.
\]

\[
(an\log^\tau n),
\qquad
\tau<0.
\]

\pause

\begin{alertblock}{Idea principal}
El método de Euler proporciona una herramienta
eficaz para estimar sumas exponenciales.
\end{alertblock}

\end{frame}
\section{Método de Van der Corput}
\begin{frame}{Desigualdad de Van der Corput}

\begin{block}{Lema}
Sean $H,N\in\mathbb N$ con $1\le H\le N$. Entonces

\[
\left|
\sum_{n=1}^{N}
e^{2\pi i f(n)}
\right|^2
\le
\frac{4N}{H}
\sum_{h=0}^{H-1}
\left|
\sum_{n=1}^{N-h}
e^{2\pi i(f(n+h)-f(n))}
\right|.
\]
\end{block}


\begin{alertblock}{Idea}
La suma exponencial asociada a $f(n)$ puede controlarse
mediante sumas exponenciales asociadas a las diferencias

\[
f(n+h)-f(n).
\]
\end{alertblock}

\end{frame}


\begin{frame}{Primer ejemplo: la sucesión $(n^2\alpha)$}

Sea

\[
f(n)=kn^2\alpha.
\]

Entonces

\[
f(n+h)-f(n)
=
k(2hn\alpha+h^2\alpha).
\]


Como $(n\alpha)$ es equidistribuida para
$\alpha\notin\mathbb Q$,

\[
(2hn\alpha+h^2\alpha)
\]

también lo es para cada $h\neq 0$.



Aplicando la desigualdad de Van der Corput y el criterio de Weyl,

\[
(n^2\alpha)
\]

es equidistribuida módulo $1$.

\end{frame}

\begin{frame}{Primer ejemplo: la sucesión $(n^2\alpha)$}

Sea

\[
f(n)=kn^2\alpha.
\]

Entonces

\[
f(n+h)-f(n)
=
k(2hn\alpha+h^2\alpha).
\]

Como $(n\alpha)$ es equidistribuida para
$\alpha\notin\mathbb Q$,

\[
(2hn\alpha+h^2\alpha)
\]

también lo es para cada $h\neq 0$.


Aplicando la desigualdad de Van der Corput y el criterio de Weyl,

\[
(n^2\alpha)
\]

es equidistribuida módulo $1$.

\end{frame}

\begin{frame}{Teorema de las diferencias de Van der Corput}

\begin{block}{Teorema}
Sea $(x_n)$ una sucesión de números reales.

Si para todo entero positivo $h$

\[
(x_{n+h}-x_n)
\]

es equidistribuida módulo $1$,

entonces

\[
(x_n)
\]

también es equidistribuida módulo $1$.
\end{block}


\begin{alertblock}{Importancia}
El resultado permite sustituir una sucesión complicada
por sus sucesiones de diferencias.

Cada aplicación reduce el orden de crecimiento
del problema.
\end{alertblock}

\[
x_n
\longrightarrow
x_{n+h}-x_n
\]

\end{frame}

\begin{frame}{Polinomios con coeficientes irracionales}

\begin{block}{Teorema}
Sea

\[
P(x)=c_mx^m+\cdots+c_1x+c_0.
\]

Si al menos uno de los coeficientes

\[
c_1,\ldots,c_m
\]

es irracional, entonces

\[
(P(n))
\]

es equidistribuida módulo $1$.
\end{block}



\begin{alertblock}{Idea de la demostración}
Las diferencias

\[
P(n+h)-P(n)
\]

son polinomios de grado menor.

Aplicando repetidamente el teorema de las diferencias,
se reduce el problema al caso lineal

\[
n\alpha,
\]

que ya conocemos por el teorema de Weyl.
\end{alertblock}

\end{frame}


\begin{frame}{Teorema de Fejér Generalizado}

\begin{block}{Teorema}
Sea $g:[1,\infty)\to\mathbb R$ una función $k$ veces derivable.

Supongamos que

\[
g^{(k)}(x)\to 0,
\]

\[
x\,|g^{(k)}(x)|\to\infty,
\]

y que $g^{(k)}$ es monótona para $x$
suficientemente grande.

Entonces

\[
(g(n))
\]

es equidistribuida módulo $1$.
\end{block}

\end{frame}

\begin{frame}{Idea de la demostración}

\[
g(n)
\]

\[
\Downarrow
\]

\[
g(n+h)-g(n)
\]

\[
\Downarrow
\]

Derivadas de orden inferior

\[
\Downarrow
\]

Teorema de Fejér
\[
(k=1)
\]


\begin{alertblock}{Idea clave}
El método de Van der Corput permite descender
desde la derivada $k$-ésima hasta el caso tratado
por el teorema de Fejér.
\end{alertblock}

\end{frame}

\begin{frame}{Aplicaciones}

\begin{block}{Consecuencias}
Son equidistribuidas módulo $1$ las sucesiones
\end{block}

\[
(an^\sigma\log^\tau n),
\qquad
0<\sigma<1,
\quad
\tau\in\mathbb R.
\]


\[
(a\log^\tau n),
\qquad
\tau>1.
\]


\[
(an\log^\tau n),
\qquad
\tau<0.
\]


\begin{alertblock}{Conclusión}
El método de Van der Corput extiende el alcance
del teorema de Fejér a familias mucho más amplias
de sucesiones.
\end{alertblock}

\end{frame}

\begin{frame}[allowframebreaks]{Bibliografía}
\nocite{*}
\printbibliography
\end{frame}

\begin{frame}
    \vspace{0.5cm}
    \begin{center}
        \Large{\bf Fin}

        \vspace{0.5cm}

        \large{Gracias por su atención}
    \end{center}
    \vspace{0.5cm}
    \begin{center}
        
\includegraphics[scale=0.29]{./pics/institucional/logos.png}
    \end{center}
\end{frame}

\end{document}